\documentclass[11pt,a4paper]{article}
\usepackage[utf8]{inputenc}
\usepackage[T1]{fontenc}
\usepackage[french]{babel}
\usepackage[a4paper,top=2cm,bottom=2.1cm,left=2.05cm,right=2.05cm]{geometry}
\usepackage{amsmath,amssymb}
\usepackage{array}
\usepackage{booktabs}
\usepackage{enumitem}
\usepackage{eurosym}
\usepackage{fancyhdr}
\usepackage{hyperref}
\usepackage{inconsolata}
\usepackage{listings}
\usepackage{microtype}
\usepackage{tabularx}
\usepackage[most]{tcolorbox}
\usepackage{titlesec}
\usepackage{xcolor}

\definecolor{ink}{HTML}{202124}
\definecolor{muted}{HTML}{5F6368}
\definecolor{line}{HTML}{DADCE0}
\definecolor{paper}{HTML}{F8F9FA}
\definecolor{accent}{HTML}{8A1538}
\definecolor{accentsoft}{HTML}{F8E8EE}
\definecolor{green}{HTML}{0B6B57}
\definecolor{greensoft}{HTML}{E6F4EF}
\definecolor{amber}{HTML}{9A5B00}
\definecolor{ambersoft}{HTML}{FFF4D6}
\definecolor{blue}{HTML}{245B8E}
\definecolor{bluesoft}{HTML}{EAF2FA}
\definecolor{violet}{HTML}{5C4D7D}
\definecolor{violetsoft}{HTML}{F1EEF8}

\hypersetup{
  colorlinks=true,
  linkcolor=accent,
  urlcolor=green
}

\setlength{\parindent}{0pt}
\setlength{\parskip}{5pt}
\setlist[itemize]{leftmargin=1.35em,itemsep=2pt,topsep=2pt}
\setlist[enumerate]{leftmargin=1.55em,itemsep=3pt,topsep=3pt}
\setlist[enumerate,1]{label=\textbf{\alph*.}}

\pagestyle{fancy}
\fancyhf{}
\fancyhead[L]{Python pour économistes}
\fancyhead[R]{Université de Lille -- L2 Économie}
\fancyfoot[C]{\thepage}
\renewcommand{\headrulewidth}{0.35pt}
\setlength{\headheight}{14.5pt}

\titleformat{\section}
  {\Large\bfseries\color{accent}}
  {}{0pt}{}
  [\vspace{-0.25em}\color{line}\titlerule]
\titleformat{\subsection}
  {\large\bfseries\color{ink}}
  {}{0pt}{}

\lstdefinelanguage{terminal}{
  morekeywords={ls,cd,pwd,mkdir,python,python3,code,dir,type,echo},
  sensitive=true
}

\lstset{
  basicstyle=\ttfamily\small,
  backgroundcolor=\color{paper},
  frame=single,
  framerule=0.35pt,
  rulecolor=\color{line},
  language=Python,
  keywordstyle=\color{accent}\bfseries,
  commentstyle=\color{muted},
  stringstyle=\color{green},
  showstringspaces=false,
  breaklines=true,
  columns=fullflexible,
  keepspaces=true,
  upquote=true,
  literate=
   {à}{{\`a}}1 {â}{{\^a}}1 {ä}{{\"a}}1
   {ç}{{\c c}}1
   {é}{{\'e}}1 {è}{{\`e}}1 {ê}{{\^e}}1 {ë}{{\"e}}1
   {î}{{\^\i}}1 {ï}{{\"\i}}1
   {ô}{{\^o}}1 {ö}{{\"o}}1
   {ù}{{\`u}}1 {û}{{\^u}}1 {ü}{{\"u}}1
   {À}{{\`A}}1 {Â}{{\^A}}1 {Ä}{{\"A}}1
   {Ç}{{\c C}}1
   {É}{{\'E}}1 {È}{{\`E}}1 {Ê}{{\^E}}1 {Ë}{{\"E}}1
   {Î}{{\^I}}1 {Ï}{{\"I}}1
   {Ô}{{\^O}}1 {Ö}{{\"O}}1
   {Ù}{{\`U}}1 {Û}{{\^U}}1 {Ü}{{\"U}}1
   {œ}{{\oe}}1 {Œ}{{\OE}}1
   {–}{{--}}1 {—}{{---}}1
}

\newcommand{\code}[1]{\texttt{#1}}
\newcommand{\pp}{\,points de pourcentage}

\newcounter{exercise}
\newcounter{extension}
\newcounter{logic}

\newcommand{\workbooktitle}[4]{%
  \setcounter{exercise}{0}%
  \setcounter{extension}{0}%
  \setcounter{logic}{0}%
  \fancyhead[L]{#1 -- #2}%
  {\Large\bfseries\color{ink}#1 -- #2\par}
  \vspace{0.2em}
  {\color{muted}Python pour économistes -- Université de Lille, L2 Économie\par}
  \vspace{0.65em}
  {\color{accent}\rule{\linewidth}{0.8pt}}
  \vspace{0.5em}
  \begin{routebox}{Question fil rouge}
  #3
  \end{routebox}
  \begin{deliverablebox}{Livrables}
  #4
  \end{deliverablebox}
  \begin{methodbox}{Méthode de travail}
  \begin{tabularx}{\linewidth}{@{}>{\bfseries\color{blue}}p{0.18\linewidth}X@{}}
  Lire & identifier les données, la question et le livrable attendu.\\
  Décomposer & écrire les étapes en pseudo-code avant de coder.\\
  Exécuter & lancer le script complet depuis le bon dossier.\\
  Vérifier & contrôler les sorties, les fichiers créés et l’interprétation.
  \end{tabularx}
  \end{methodbox}
}

\newtcolorbox{routebox}[1]{
  enhanced,breakable,
  title={#1},
  colback=paper,
  colframe=line,
  coltitle=ink,
  fonttitle=\bfseries,
  boxrule=0.5pt,
  arc=2pt,
  left=8pt,right=8pt,top=7pt,bottom=7pt,
  borderline west={2.4pt}{0pt}{accent}
}

\newtcolorbox{deliverablebox}[1]{
  enhanced,breakable,
  title={#1},
  colback=greensoft,
  colframe=green,
  coltitle=ink,
  fonttitle=\bfseries,
  boxrule=0.6pt,
  arc=2pt,
  left=8pt,right=8pt,top=7pt,bottom=7pt,
  borderline west={2.4pt}{0pt}{accent}
}


\newtcolorbox{methodbox}[1]{
  enhanced,breakable,
  title={#1},
  colback=bluesoft,
  colframe=blue,
  coltitle=ink,
  fonttitle=\bfseries,
  boxrule=0.55pt,
  arc=2pt,
  left=8pt,right=8pt,top=6pt,bottom=6pt,
  before skip=7pt,after skip=9pt,
  borderline west={2.4pt}{0pt}{blue}
}

\newtcolorbox{pathwaybox}[1]{
  enhanced,breakable,
  title={#1},
  colback=violetsoft,
  colframe=violet,
  coltitle=ink,
  fonttitle=\bfseries,
  boxrule=0.55pt,
  arc=2pt,
  left=8pt,right=8pt,top=6pt,bottom=6pt,
  before skip=8pt,after skip=8pt,
  borderline west={2.4pt}{0pt}{violet}
}

\newenvironment{pseudocodeexercise}[1]{%
  \refstepcounter{logic}%
  \begin{tcolorbox}[
    enhanced,breakable,
    title={Logique \thelogic{} -- #1 \hfill\scriptsize PSEUDO-CODE},
    colback=bluesoft,
    colframe=blue,
    coltitle=ink,
    fonttitle=\bfseries,
    boxrule=0.6pt,
    arc=2pt,
    left=8pt,right=8pt,top=6pt,bottom=6pt,
    before skip=8pt,after skip=8pt,
    borderline west={2.6pt}{0pt}{blue}
  ]%
}{\end{tcolorbox}}

\newenvironment{mandatoryexercise}[1]{%
  \refstepcounter{exercise}%
  \begin{tcolorbox}[
    enhanced,breakable,
    title={Exercice \theexercise{} -- #1 \hfill\scriptsize OBLIGATOIRE},
    colback=white,
    colframe=accent,
    coltitle=ink,
    fonttitle=\bfseries,
    boxrule=0.7pt,
    arc=2pt,
    left=8pt,right=8pt,top=7pt,bottom=7pt,
    before skip=9pt,after skip=9pt,
    borderline west={2.4pt}{0pt}{accent}
  ]%
}{\end{tcolorbox}}

\newenvironment{extensionexercise}[1]{%
  \refstepcounter{extension}%
  \begin{tcolorbox}[
    enhanced,breakable,
    title={Extension \theextension{} -- #1},
    colback=ambersoft,
    colframe=amber,
    coltitle=ink,
    fonttitle=\bfseries,
    boxrule=0.55pt,
    arc=2pt,
    left=8pt,right=8pt,top=7pt,bottom=7pt,
    before skip=8pt,after skip=8pt,
    borderline west={2.4pt}{0pt}{amber}
  ]%
}{\end{tcolorbox}}

\newtcolorbox{checkpointbox}[1]{
  enhanced,breakable,
  title={#1},
  colback=accentsoft,
  colframe=accent,
  coltitle=ink,
  fonttitle=\bfseries,
  boxrule=0.55pt,
  arc=2pt,
  left=8pt,right=8pt,top=7pt,bottom=7pt,
  borderline west={2.4pt}{0pt}{accent}
}


\begin{document}

\workbooktitle
  {TP2}
  {Boucles, conditions et proportions avec données OWID}
  {Comment utiliser des boucles et des conditions pour classer des observations économiques simples ?}
  {Un fichier \code{tp2\_nom\_prenom.py} contenant une boucle, une condition, une proportion et une interprétation économique courte.}

\begin{checkpointbox}{Parcours obligatoire}
Les exercices 1 à 7 sont obligatoires. La notion statistique principale est la proportion. Les seuils utilisés sont pédagogiques : ils servent à apprendre \code{if}, \code{for} et les compteurs.
\end{checkpointbox}

\begin{checkpointbox}{Source des données}
Les valeurs de CO$_2$ par habitant sont des extraits pédagogiques arrondis issus d'Our World in Data.
\begin{lstlisting}[language=terminal]
https://ourworldindata.org/grapher/co-emissions-per-capita.csv
\end{lstlisting}
\end{checkpointbox}

\begin{pseudocodeexercise}{Classer des observations}
\begin{lstlisting}
# Objectif : classer des pays selon un indicateur
# Entree : liste de pays et valeurs de CO2 par habitant
# Pour chaque pays :
#     lire la valeur
#     comparer la valeur a un seuil
#     afficher une categorie
# Compter combien de pays sont dans chaque categorie
# Sortie : effectifs, proportions, interpretation
\end{lstlisting}
\end{pseudocodeexercise}

\section*{A. Représenter les données}

\begin{mandatoryexercise}{Créer une liste de dictionnaires}
Créer la liste suivante. Chaque dictionnaire représente une observation.
\begin{lstlisting}
donnees = [
    {"pays": "France", "co2_pc": 4.7},
    {"pays": "Germany", "co2_pc": 7.9},
    {"pays": "Spain", "co2_pc": 5.0},
    {"pays": "Poland", "co2_pc": 8.4},
    {"pays": "United States", "co2_pc": 14.9},
    {"pays": "India", "co2_pc": 2.0},
    {"pays": "China", "co2_pc": 8.0},
    {"pays": "World", "co2_pc": 4.7}
]
\end{lstlisting}
Afficher le nombre d'observations et la première observation.
\end{mandatoryexercise}

\begin{mandatoryexercise}{Parcourir les observations}
Avec une boucle \code{for}, afficher une phrase par pays :
\begin{lstlisting}
for obs in donnees:
    print(f"{obs['pays']} : {obs['co2_pc']} tonnes par habitant")
\end{lstlisting}
Modifier ensuite l'affichage pour arrondir à une décimale.
\end{mandatoryexercise}

\section*{B. Conditions et proportions}

\begin{mandatoryexercise}{Classer avec une condition}
On utilise trois catégories pédagogiques :
\begin{itemize}
  \item \code{faible} si \code{co2\_pc < 5} ;
  \item \code{intermediaire} si \code{5 <= co2\_pc < 10} ;
  \item \code{eleve} si \code{co2\_pc >= 10}.
\end{itemize}
Écrire une boucle qui affiche la catégorie de chaque pays.
\end{mandatoryexercise}

\begin{mandatoryexercise}{Compter une proportion}
Compter le nombre de pays dont les émissions dépassent 5 tonnes par habitant. Calculer ensuite la proportion :
\[
\frac{\text{nombre de pays au-dessus du seuil}}{\text{nombre total de pays}} \times 100.
\]
Afficher une phrase du type : \code{62.5 \% des observations sont au-dessus du seuil.}
\end{mandatoryexercise}

\begin{mandatoryexercise}{Créer une fonction de classement}
Écrire une fonction \code{classe\_co2(valeur)} qui renvoie \code{"faible"}, \code{"intermediaire"} ou \code{"eleve"}. L'utiliser dans une boucle pour afficher le classement de tous les pays.
\end{mandatoryexercise}

\section*{C. Interpréter}

\begin{mandatoryexercise}{Question économique}
Ajouter quatre commentaires à la fin du script :
\begin{enumerate}
  \item quel pays de l'extrait a la valeur la plus élevée ?
  \item la France est-elle au-dessus ou au-dessous de la valeur mondiale ?
  \item quelle proportion de l'extrait dépasse 5 tonnes par habitant ?
  \item pourquoi ne faut-il pas confondre CO$_2$ par habitant et CO$_2$ total ?
\end{enumerate}
\end{mandatoryexercise}

\begin{mandatoryexercise}{Contrôle final}
Le script doit contenir au minimum : une liste de dictionnaires, une boucle \code{for}, une condition \code{if/elif/else}, un compteur, une proportion, une fonction et une interprétation.
\end{mandatoryexercise}

\section*{D. Entraînement facultatif}

\begin{extensionexercise}{Dix petits entraînements}
Choisir au moins deux questions si le parcours obligatoire est terminé.
\begin{enumerate}
  \item Changer le seuil de 5 à 7 tonnes et recalculer la proportion.
  \item Créer une liste contenant seulement les pays classés \code{eleve}.
  \item Calculer la moyenne de \code{co2\_pc} avec une boucle, sans \code{sum()}.
  \item Ajouter un pays et vérifier que le code fonctionne encore.
  \item Afficher les pays au-dessous de la valeur mondiale.
  \item Compter séparément les pays \code{faible}, \code{intermediaire} et \code{eleve}.
  \item Écrire une fonction \code{au\_dessus\_seuil(valeur, seuil)}.
  \item Trier manuellement deux pays : lequel a la valeur la plus élevée ?
  \item Écrire une phrase sur la limite d'un classement par seuil.
  \item Expliquer pourquoi une proportion dépend de l'échantillon choisi.
\end{enumerate}
\end{extensionexercise}

\end{document}
